\begin{figure}[ht]
\centering
\begin{tikzpicture}[x=1.0cm, y=1.0cm, font=\small,
  node distance=2.2cm,
  vtx/.style={circle, draw=engblack, thick, minimum size=8mm,
    inner sep=1pt, font=\scriptsize},
  edgelbl/.style={font=\scriptsize, fill=white, inner sep=1pt}]

% Nodes: source s, two intermediate a,b, sink t. A small diamond network.
\node[vtx, fill=englime!30] (s) at (0, 0) {$s$};
\node[vtx] (a) at (2.8, 1.4) {$a$};
\node[vtx] (b) at (2.8, -1.4) {$b$};
\node[vtx, fill=engorange!30] (t) at (5.6, 0) {$t$};

% Edges with (cost, capacity) labels and the optimal flow drawn thick.
% s->a: cost 2, cap 3, flow 3 (saturated).
\draw[engblue, very thick, -{Latex[length=2mm]}] (s) -- (a)
  node[edgelbl, midway, above left] {$2,\,3$};
% s->b: cost 1, cap 2, flow 2 (saturated).
\draw[engblue, very thick, -{Latex[length=2mm]}] (s) -- (b)
  node[edgelbl, midway, below left] {$1,\,2$};
% a->t: cost 1, cap 4, flow 3.
\draw[engblue, thick, -{Latex[length=2mm]}] (a) -- (t)
  node[edgelbl, midway, above right] {$1,\,4$};
% b->t: cost 3, cap 2, flow 2 (saturated).
\draw[engblue, very thick, -{Latex[length=2mm]}] (b) -- (t)
  node[edgelbl, midway, below right] {$3,\,2$};
% a->b: cost 1, cap 2, flow 0 (idle cross-link).
\draw[enggray!70, thick, dashed, -{Latex[length=2mm]}] (a) -- (b)
  node[edgelbl, midway] {$1,\,2$};

% Supply/demand annotations.
\node[englime, anchor=east, font=\scriptsize] at (-0.55, 0) {$+5$};
\node[engorange, anchor=west, font=\scriptsize] at (6.15, 0) {$-5$};

\node[anchor=north, font=\scriptsize] at (2.8, -2.4)
  {edge label $=$ (cost per unit, capacity); thick $=$ carrying flow};

\end{tikzpicture}
\caption[A minimum-cost network-flow instance]{A minimum-cost flow
routing five units from the source $s$ to the sink $t$ across a small
capacitated network. Each edge carries a per-unit cost and a capacity,
shown as the label pair, and the decision is how much flow to send on
each edge to meet the demand at least total cost while respecting every
capacity and conserving flow at each intermediate node. The optimum
saturates the two cheap source edges and the expensive $b\!\to\!t$ edge
and leaves the cross-link $a\!\to\!b$ idle (dashed), a pattern the
active set reads off directly. The constraint matrix of a flow problem
is the node-edge incidence matrix, which is totally unimodular, so the
LP relaxation returns integer flows with no branching: a network flow
is an integer program that the polynomial-time LP solver already
solves. The dual variables are node potentials, and their differences
along an edge are the reduced costs that a residual-network algorithm
drives to optimality, the same shadow-price reading that priced machine
hours in the production LP now pricing congestion at each node.}
\label{fig:vol02:ch15:network-flow}
\end{figure}
